\documentclass[preview]{standalone}

\usepackage{amsmath}
\usepackage{amssymb}
\usepackage{stellar}
\usepackage{definitions}
\usepackage{bettelini}
\usepackage{xfrac}

\begin{document}

\id{lebesgue-measure-exercises2}
\genpage

\section{Exercises}

\begin{snippetsolution}{lebesgue-measure-ex-12-sol}{}
    Let
    \[
        I_n
        =
        \integral[0][1][
            (1+nx^2)(1+x^2)^{-n}
        ][x].
    \]
    For every fixed \(x\in(0,1]\), we have
    \[
        (1+x^2)^{-n}\to 0,
    \]
    and therefore
    \[
        (1+nx^2)(1+x^2)^{-n}\to 0.
    \]
    The point \(x=0\) is negligible, so the pointwise limit is \(0\) almost
    everywhere.

    However, instead of applying the dominated convergence theorem directly, we
    prove a uniform estimate. Since \(0\leq x\leq 1\), we have
    \[
        \log(1+x^2)\geq \frac{x^2}{2}.
    \]
    Hence
    \[
        (1+x^2)^{-n}
        =
        e^{-n\log(1+x^2)}
        \leq
        e^{-\sfrac{nx^2}{2}}.
    \]
    Therefore
    \[
        0\leq I_n
        \leq
        \integral[0][1][
            (1+nx^2)e^{-\sfrac{nx^2}{2}}
        ][x].
    \]
    With the change of variables
    \[
        y=\sqrt n\,x,
    \]
    we get
    \[
        I_n
        \leq
        \frac1{\sqrt n}
        \integral[0][\sqrt n][
            (1+y^2)e^{-\sfrac{y^2}{2}}
        ][y].
    \]
    Since
    \[
        \integral[0][+\infty][
            (1+y^2)e^{-\sfrac{y^2}{2}}
        ][y]<+\infty,
    \]
    it follows that
    \[
        0\leq I_n\leq \frac{C}{\sqrt n}
    \]
    for some constant \(C>0\). Hence
    \[
        \lim_{n\to\infty} I_n=0.
    \]

    Therefore
    \[
        \lim_{n\to\infty}
        \integral[0][1][
            (1+nx^2)(1+x^2)^{-n}
        ][x]
        =
        0
    \]
\end{snippetsolution}

\begin{snippetexercise}{lebesgue-measure-ex-13}{}
    Compute \[
        \lim_n \integral[0][+\infty][
            n \sin \left(\frac xn\right) \left[x(1+x^2)\right]^{-1}
        ][x]
    \]
\end{snippetexercise}

\begin{snippetsolution}{lebesgue-measure-ex-13-sol}{}
    Let
    \[
        f_n(x)
        =
        n\sin\left(\frac xn\right)\left[x(1+x^2)\right]^{-1},
        \qquad x>0.
    \]
    For every fixed \(x>0\), we have
    \[
        n\sin\left(\frac xn\right)
        =
        x\frac{\sin(\sfrac xn)}{\sfrac xn}
        \to x.
    \]
    Therefore
    \[
        f_n(x)\to \frac{1}{1+x^2}.
    \]

    Moreover, since \(0\leq \sin u\leq u\) for \(u\geq 0\), we have
    \[
        0\leq
        n\sin\left(\frac xn\right)
        \leq x.
    \]
    Hence
    \[
        0\leq f_n(x)\leq \frac{1}{1+x^2}.
    \]
    Since
    \[
        \frac{1}{1+x^2}\in\mathcal L^1((0,+\infty)),
    \]
    the dominated convergence theorem gives
    \[
        \lim_{n\to\infty}
        \integral[0][+\infty][
            n \sin \left(\frac xn\right) \left[x(1+x^2)\right]^{-1}
        ][x]
        =
        \integral[0][+\infty][\frac{1}{1+x^2}][x].
    \]
    Therefore
    \[
        \lim_{n\to\infty}
        \integral[0][+\infty][
            n \sin \left(\frac xn\right) \left[x(1+x^2)\right]^{-1}
        ][x]
        =
        \frac\pi2
    \]
\end{snippetsolution}

\begin{snippetexercise}{lebesgue-measure-ex-14}{}
    Compute \[
        \lim_n \integral[a][\infty][n(1 + n^2x^2)^{-1}][x]
    \]
    \emph{Hint:} the answers depends on whether \(a>0, a=0, a<0\).
\end{snippetexercise}

\begin{snippetsolution}{lebesgue-measure-ex-14-sol}{}
    Let
    \[
        I_n(a)
        =
        \integral[a][\infty][\frac{n}{1+n^2x^2}][x].
    \]
    We use the change of variables
    \[
        u=nx,
        \qquad
        du=n\,dx.
    \]
    Then
    \[
        I_n(a)
        =
        \integral[na][+\infty][\frac{1}{1+u^2}][u].
    \]
    Therefore
    \[
        I_n(a)
        =
        \left[\arctan u\right]_{na}^{+\infty}
        =
        \frac\pi2-\arctan(na).
    \]

    Now we distinguish three cases.

    If \(a>0\), then \(na\to+\infty\), hence
    \[
        \arctan(na)\to \frac\pi2.
    \]
    Therefore
    \[
        \lim_{n\to\infty} I_n(a)=0.
    \]

    If \(a=0\), then
    \[
        I_n(0)
        =
        \integral[0][+\infty][\frac{n}{1+n^2x^2}][x]
        =
        \frac\pi2
    \]
    for every \(n\). Hence
    \[
        \lim_{n\to\infty} I_n(0)=\frac\pi2.
    \]

    If \(a<0\), then \(na\to-\infty\), hence
    \[
        \arctan(na)\to -\frac\pi2.
    \]
    Therefore
    \[
        \lim_{n\to\infty} I_n(a)
        =
        \frac\pi2-\left(-\frac\pi2\right)
        =
        \pi.
    \]

    Consequently,
    \[
        \lim_{n\to\infty}
        \integral[a][\infty][n(1+n^2x^2)^{-1}][x]
        =
        \begin{cases}
            0 & a>0,\\
            \frac\pi2 & a=0,\\
            \pi & a<0.
        \end{cases}
    \]
\end{snippetsolution}

\begin{snippetexercise}{lebesgue-measure-ex-15}{}
    Study the \function \[
        F(x) = \integral[0][x][\frac{\ln\left(1 + \frac tx\right)}{1 + t^2}][t]
    \]
\end{snippetexercise}

\begin{snippetsolution}{lebesgue-measure-ex-15-sol}{}
    The function is naturally defined for \(x\neq 0\). Indeed, if \(x>0\),
    then \(t\in[0,x]\) implies
    \[
        1+\frac tx\in[1,2],
    \]
    and the same is true if \(x<0\), taking into account the orientation of the
    integral. Thus the natural domain is
    \[
        \realnumbers\setminus\{0\}.
    \]

    For \(x\neq 0\), use the change of variables
    \[
        t=xs.
    \]
    Since \(t\) goes from \(0\) to \(x\), the variable \(s\) goes from \(0\) to
    \(1\). Hence
    \[
        F(x)
        =
        x\integral[0][1][
            \frac{\ln(1+s)}{1+x^2s^2}
        ][s].
    \]
    This formula shows immediately that \(F\) is odd:
    \[
        F(-x)=-F(x).
    \]

    Moreover,
    \[
        |F(x)|
        \leq
        |x|\integral[0][1][\ln(1+s)][s],
    \]
    so
    \[
        \lim_{x\to 0}F(x)=0.
    \]
    Therefore the singularity at \(x=0\) is removable. If we define
    \[
        F(0)=0,
    \]
    then \(F\) becomes continuous on the whole real line.

    We now compute the derivative at \(0\). For \(x\neq 0\),
    \[
        \frac{F(x)-F(0)}{x}
        =
        \integral[0][1][
            \frac{\ln(1+s)}{1+x^2s^2}
        ][s].
    \]
    Since
    \[
        0\leq
        \frac{\ln(1+s)}{1+x^2s^2}
        \leq
        \ln(1+s)
        \in\mathcal L^1(0,1),
    \]
    by the dominated convergence theorem,
    \[
        F'(0)
        =
        \integral[0][1][\ln(1+s)][s].
    \]
    Hence
    \[
        F'(0)
        =
        \integral[1][2][\ln u][u]
        =
        \left[u\ln u-u\right]_1^2
        =
        2\ln2-1.
    \]

    For \(x\neq0\), differentiating under the integral sign is justified on
    every compact subset of \(\realnumbers\), because the integrand and its
    derivative with respect to \(x\) are dominated by integrable functions on
    \((0,1)\). We get
    \[
        F'(x)
        =
        \integral[0][1][
            \ln(1+s)\frac{1-x^2s^2}{(1+x^2s^2)^2}
        ][s].
    \]
    In order to study the sign of \(F'\), it is convenient to compute it more
    explicitly. For \(x>0\), by Leibniz' rule,
    \[
        F'(x)
        =
        \frac{\ln2}{1+x^2}
        -
        \frac1x
        \integral[0][x][
            \frac{t}{(x+t)(1+t^2)}
        ][t].
    \]
    A partial fraction decomposition gives
    \[
        \frac{t}{(x+t)(1+t^2)}
        =
        -\frac{x}{1+x^2}\frac1{x+t}
        +
        \frac{x}{1+x^2}\frac{t}{1+t^2}
        +
        \frac1{1+x^2}\frac1{1+t^2}.
    \]
    Therefore
    \[
        \integral[0][x][
            \frac{t}{(x+t)(1+t^2)}
        ][t]
        =
        \frac{
            -x\ln2+\frac x2\ln(1+x^2)+\arctan x
        }{1+x^2}.
    \]
    Hence, for \(x>0\),
    \[
        F'(x)
        =
        \frac{
            2\ln2-\frac12\ln(1+x^2)-\frac{\arctan x}{x}
        }{1+x^2}.
    \]
    Since \(F\) is odd, \(F'\) is even. Thus the same formula holds for
    \(x\neq0\).

    Define, for \(x>0\),
    \[
        h(x)
        =
        2\ln2-\frac12\ln(1+x^2)-\frac{\arctan x}{x}.
    \]
    Then
    \[
        h'(x)
        =
        \frac{\arctan x-x}{x^2}<0.
    \]
    Moreover,
    \[
        \lim_{x\to0^+}h(x)=2\ln2-1>0,
        \qquad
        \lim_{x\to+\infty}h(x)=-\infty.
    \]
    Hence there exists a unique \(\alpha>0\) such that
    \[
        h(\alpha)=0.
    \]
    Equivalently,
    \[
        2\ln2-\frac12\ln(1+\alpha^2)-\frac{\arctan\alpha}{\alpha}=0.
    \]
    Numerically,
    \[
        \alpha\simeq 2.1263.
    \]

    Therefore \(F\) is decreasing on \((-\infty,-\alpha)\), increasing on
    \((-\alpha,\alpha)\), and decreasing on \((\alpha,+\infty)\). Consequently,
    \(F\) has a global minimum at \(-\alpha\) and a global maximum at
    \(\alpha\).

    Finally, for \(x>0\),
    \[
        0\leq F(x)
        =
        x\integral[0][1][
            \frac{\ln(1+s)}{1+x^2s^2}
        ][s]
        \leq
        x\integral[0][1][
            \frac{s}{1+x^2s^2}
        ][s]
        =
        \frac{\ln(1+x^2)}{2x}.
    \]
    Hence
    \[
        \lim_{x\to+\infty}F(x)=0.
    \]
    Since \(F\) is odd,
    \[
        \lim_{x\to-\infty}F(x)=0.
    \]

    Thus the graph is odd, has a removable zero at the origin, increases through
    the origin with slope
    \[
        F'(0)=2\ln2-1,
    \]
    reaches a positive maximum at \(x=\alpha\), then decreases to the horizontal
    asymptote \(y=0\). By oddness, it has a symmetric negative minimum at
    \(x=-\alpha\).
\end{snippetsolution}

\begin{snippetexercise}{lebesgue-measure-ex-16}{}
    Let \[
        S = \{
            (x,y,z) \in \realnumbers^3 \suchthat
            \sqrt{x^2 + y^2} < z < \sqrt{1 - x^2 - y^2}
        \}
    \]
    and consider the \sequence of \function[functions]
    \[
        f_n(x,y,z) = (\sinh y)\left(
            x + \frac 1n
        \right)^{\sfrac13} e^{\frac yn - 3\ln z}
    \]
    Show that \(f_n \in \mathcal L^1(S)\) for each \(n\) and compute \[
        \lim_n \int_S f_n\,d\mu
    \]
\end{snippetexercise}

\begin{snippetsolution}{lebesgue-measure-ex-16-sol}{}
    We first rewrite the set \(S\) in cylindrical coordinates:
    \[
        x=r\cos\theta,\qquad y=r\sin\theta.
    \]
    The condition defining \(S\) becomes
    \[
        r<z<\sqrt{1-r^2}.
    \]
    This is possible if and only if
    \[
        r<\frac1{\sqrt2}.
    \]
    Hence
    \[
        S=
        \left\{
            (r,\theta,z)\suchthat
            0\leq r<\frac1{\sqrt2},\,
            0\leq \theta<2\pi,\,
            r<z<\sqrt{1-r^2}
        \right\}.
    \]
    In particular, \(S\) is bounded and \(z>0\) on \(S\).

    Since \(S\) is bounded, there exists \(C>0\), independent of \(n\), such that
    \[
        \left|x+\frac1n\right|^{\sfrac13}e^{\sfrac yn}
        \leq C
    \]
    on \(S\). Moreover, since \(y\) is bounded on \(S\),
    \[
        |\sinh y|\leq C|y|.
    \]
    Therefore
    \[
        |f_n(x,y,z)|
        \leq
        C\frac{|y|}{z^3}.
    \]
    We prove that the right-hand side is integrable on \(S\). In cylindrical
    coordinates,
    \[
        \integral[S]{}{\frac{|y|}{z^3}}[(x,y,z)]
        =
        \integral[0][2\pi]{|\sin\theta|}[\theta]
        \integral[0][\sfrac1{\sqrt2}]{
            r^2\integral[r][\sqrt{1-r^2}]{\frac1{z^3}}[z]
        }[r].
    \]
    For \(r>0\),
    \[
        \integral[r][\sqrt{1-r^2}]{\frac1{z^3}}[z]
        \leq
        \integral[r][+\infty]{\frac1{z^3}}[z]
        =
        \frac{1}{2r^2}.
    \]
    Hence
    \[
        r^2\integral[r][\sqrt{1-r^2}]{\frac1{z^3}}[z]
        \leq
        \frac12.
    \]
    Thus
    \[
        \integral[S]{}{\frac{|y|}{z^3}}[(x,y,z)]<+\infty.
    \]
    It follows that
    \[
        f_n\in\mathcal L^1(S)
    \]
    for every \(n\).

    Now, for every \((x,y,z)\in S\),
    \[
        \left(x+\frac1n\right)^{\sfrac13}\to x^{\sfrac13},
        \qquad
        e^{\sfrac yn}\to 1.
    \]
    Therefore
    \[
        f_n(x,y,z)
        \to
        f(x,y,z)
        :=
        (\sinh y)x^{\sfrac13}z^{-3}.
    \]
    The previous estimate gives an integrable dominating function independent of
    \(n\). Hence, by the dominated convergence theorem,
    \[
        \lim_{n\to\infty}\int_S f_n\,d\mu
        =
        \int_S
        (\sinh y)x^{\sfrac13}z^{-3}\,d\mu.
    \]

    The set \(S\) is symmetric with respect to the transformation
    \[
        (x,y,z)\mapsto (-x,y,z).
    \]
    Moreover,
    \[
        (\sinh y)(-x)^{\sfrac13}z^{-3}
        =
        -(\sinh y)x^{\sfrac13}z^{-3}.
    \]
    Hence the limit integrand is odd with respect to \(x\), and it is integrable
    on \(S\). Therefore its integral over \(S\) is zero. Consequently,
    \[
        \lim_{n\to\infty}\int_S f_n\,d\mu=0
    \]
\end{snippetsolution}

\begin{snippetexercise}{lebesgue-measure-ex-17}{}
    Let \[
        f_n(x,y) = \frac{\cos^2\left(\frac xn\right) e^{-\sfrac{y^2}{n}}}{
            1 + x^2 + y^2
        }
    \]
    Determine the poinwise limite of the \sequence and then determine whether:
    \begin{enumerate}
        \item the convergence of \(f_n\) is monotone;
        \item the \(f_n\) are integrable on \(\realnumbers^2\) for each \(n\). If yes stufy the limit on \(\int_{\realnumbers^2} f_n\,d\mu\).
        \item there exists a \function which bounds \(f_n\) from above for each \(n\) and such that such \function is in \(\mathcal L^1(\realnumbers^2)\).
    \end{enumerate}
\end{snippetexercise}

\begin{snippetsolution}{lebesgue-measure-ex-17-sol}{}
    For every \((x,y)\in\realnumbers^2\), we have
    \[
        \cos^2\left(\frac xn\right)\to 1,
        \qquad
        e^{-\sfrac{y^2}{n}}\to 1.
    \]
    Hence
    \[
        f_n(x,y)\to f(x,y):=\frac{1}{1+x^2+y^2}.
    \]

    The convergence is not monotone. Indeed, at the point \((0,1)\),
    \[
        f_n(0,1)=\frac{e^{-\sfrac1n}}{2},
    \]
    which is increasing in \(n\). On the other hand, at the point \((\pi,0)\),
    \[
        f_1(\pi,0)=\frac{1}{1+\pi^2},
        \qquad
        f_2(\pi,0)=0.
    \]
    Thus the sequence is not increasing. Since it increases at \((0,1)\), it is
    not decreasing either. Therefore the convergence is not monotone.

    We now prove that \(f_n\in\mathcal L^1(\realnumbers^2)\) for every \(n\).
    Since
    \[
        0\leq \cos^2\left(\frac xn\right)\leq 1,
    \]
    we have
    \[
        0\leq f_n(x,y)
        \leq
        \frac{e^{-\sfrac{y^2}{n}}}{1+x^2+y^2}
        \leq
        \frac{e^{-\sfrac{y^2}{n}}}{1+x^2}.
    \]
    Therefore
    \[
        \int_{\realnumbers^2} f_n\,d\mu
        \leq
        \int_{\realnumbers}e^{-\sfrac{y^2}{n}}\,dy
        \int_{\realnumbers}\frac{1}{1+x^2}\,dx
        <+\infty.
    \]
    Hence
    \[
        f_n\in\mathcal L^1(\realnumbers^2)
    \]
    for each \(n\).

    However, the limit function is not integrable on \(\realnumbers^2\). Indeed,
    using polar coordinates,
    \[
        \int_{\realnumbers^2}\frac{1}{1+x^2+y^2}\,d\mu
        =
        2\pi\int_0^{+\infty}\frac{r}{1+r^2}\,dr
        =
        +\infty.
    \]
    Since \(f_n\geq 0\), Fatou's lemma gives
    \[
        \int_{\realnumbers^2} f\,d\mu
        \leq
        \liminf_{n\to\infty}
        \int_{\realnumbers^2} f_n\,d\mu.
    \]
    The left-hand side is \(+\infty\), hence
    \[
        \lim_{n\to\infty}
        \int_{\realnumbers^2} f_n\,d\mu
        =
        +\infty.
    \]

    Finally, there cannot exist a function
    \[
        g\in\mathcal L^1(\realnumbers^2)
    \]
    such that
    \[
        0\leq f_n\leq g
        \qquad
        \text{for every }n.
    \]
    Indeed, if such a \(g\) existed, then by the dominated convergence theorem we
    would obtain
    \[
        \int_{\realnumbers^2} f_n\,d\mu
        \to
        \int_{\realnumbers^2} f\,d\mu.
    \]
    But this is impossible, because \(f\notin\mathcal L^1(\realnumbers^2)\).
    Equivalently, the pointwise limit of a dominated sequence would have to be
    integrable, while here
    \[
        f(x,y)=\frac{1}{1+x^2+y^2}\notin\mathcal L^1(\realnumbers^2).
    \]
\end{snippetsolution}

\begin{snippetexercise}{lebesgue-measure-ex-18}{}
    Verify that the \function \[
        f(x,y,z) = \left(
            \frac{x}{y^2 - 2^{-\sfrac12}|x|}
        \right)^{\sfrac13} \frac{1}{y(1 + z^2)}
    \]
    is integrable on \[
        D = \left\{
            (x,y,z) \in \realnumbers^3 \suchthat z \leq x^2 + y^2 \leq 1  \land |y| > |x|
        \right\}
    \]
    and compute \[
        \int_D f\,d\mu
    \]
\end{snippetexercise}

\begin{snippetsolution}{lebesgue-measure-ex-18-sol}{}
    The possible singular set
    \[
        \left\{
            (x,y,z)\in D\suchthat
            y^2=2^{-\sfrac12}|x|
        \right\}
    \]
    is negligible, hence it is irrelevant for Lebesgue integrability.

    We use cylindrical coordinates
    \[
        x=r\cos\theta,\qquad y=r\sin\theta.
    \]
    The condition \(|y|>|x|\) becomes
    \[
        |\sin\theta|>|\cos\theta|,
    \]
    and \(x^2+y^2\leq 1\) becomes \(0\leq r\leq1\). Moreover,
    \[
        z\leq x^2+y^2=r^2.
    \]
    Hence
    \[
        D=
        \left\{
            (r,\theta,z)\suchthat
            0\leq r\leq1,\,
            |\sin\theta|>|\cos\theta|,\,
            z\leq r^2
        \right\}.
    \]

    We first prove absolute integrability. Since
    \[
        \integral[-\infty][r^2][\frac{1}{1+z^2}][z]
        =
        \frac\pi2+\arctan(r^2)
        \leq \pi,
    \]
    it is enough to study the integral in the variables \(r,\theta\). We have
    \[
        y^2-2^{-\sfrac12}|x|
        =
        r^2\sin^2\theta
        -
        2^{-\sfrac12}r|\cos\theta|
        =
        r\left(
            r\sin^2\theta
            -
            2^{-\sfrac12}|\cos\theta|
        \right).
    \]
    Therefore
    \[
        \left|
            \left(
                \frac{x}{y^2-2^{-\sfrac12}|x|}
            \right)^{\sfrac13}
            \frac1y
        \right|r
        =
        \frac{
            |\cos\theta|^{\sfrac13}
        }{
            |\sin\theta|
            \left|
                r\sin^2\theta
                -
                2^{-\sfrac12}|\cos\theta|
            \right|^{\sfrac13}
        }.
    \]
    Since \(|\sin\theta|>|\cos\theta|\), we have
    \[
        |\sin\theta|\geq \frac1{\sqrt2}
    \]
    on the angular region under consideration. Thus the factor
    \(|\sin\theta|^{-1}\) is bounded.

    For fixed \(\theta\), write
    \[
        r\sin^2\theta
        -
        2^{-\sfrac12}|\cos\theta|
        =
        \sin^2\theta
        \left(
            r-r_\theta
        \right),
        \qquad
        r_\theta=
        \frac{|\cos\theta|}{\sqrt2\sin^2\theta}.
    \]
    Since \(|\sin\theta|>|\cos\theta|\), one has \(0\leq r_\theta<1\). Hence
    \[
        \integral[0][1][
            \frac{1}{
                \left|
                    r\sin^2\theta
                    -
                    2^{-\sfrac12}|\cos\theta|
                \right|^{\sfrac13}
            }
        ][r]
        \leq
        C\integral[0][1][\frac1{|r-r_\theta|^{\sfrac13}}][r]
        \leq C.
    \]
    Consequently,
    \[
        \int_D |f|\,d\mu
        \leq
        C
        \integral[|\sin\theta|>|\cos\theta|]{}{
            |\cos\theta|^{\sfrac13}
        }[\theta]
        <+\infty.
    \]
    Therefore
    \[
        f\in\mathcal L^1(D).
    \]

    It remains to compute the integral. The domain \(D\) is symmetric with
    respect to the transformation
    \[
        (x,y,z)\mapsto (-x,y,z).
    \]
    Moreover,
    \[
        f(-x,y,z)
        =
        \left(
            \frac{-x}{y^2-2^{-\sfrac12}|x|}
        \right)^{\sfrac13}
        \frac{1}{y(1+z^2)}
        =
        -f(x,y,z).
    \]
    Since \(f\in\mathcal L^1(D)\), the symmetry argument is justified, and we
    obtain
    \[
        \int_D f\,d\mu=0.
    \]

    Hence
    \[
        \int_D f\,d\mu=0
    \]
\end{snippetsolution}

\begin{snippetexercise}{lebesgue-measure-ex-19}{}
    Let \begin{align*}
        f_n(x,y) &= \frac{y}{1 + x^4 + y^4} \operatorname{exp}\left(
            \frac{1}{x^3 - n|y|}
        \right) \\
        A_n &= \left\{
            (x,y) \in \realnumbers^2 \suchthat x^3 - n|y| < 0
        \right\} \\
        f &= \sum_{n=1}^\infty \frac{1}{n^2} f_n \cdot \chi_{A_n}
    \end{align*}
    \begin{enumerate}
        \item Show that \(f\) is defined almost everywhere;
        \item Show that \(f\) is non-null almost everywhere;
        \item Show that \(f\in \mathcal L^1(\realnumbers^2)\) and compute \(\int_{\realnumbers^2} f\,d\mu\);
        \item Show that, for each \(n = 1, 2, \cdots\), \(f_n \notin \mathcal L^1(\realnumbers^2 \difference A_n)\);
        \item If \(h(z)\) is a real \function of real variables, find a sufficient and necessary
        condition on \(h\) such that \(h(z) \cdot f(x,y)\) is in \(\mathcal L^1(\realnumbers^3)\).
    \end{enumerate}
\end{snippetexercise}

\begin{snippetsolution}{lebesgue-measure-ex-19-sol}{}
    Let
    \[
        E_n=\left\{(x,y)\in\realnumbers^2\suchthat x^3-n|y|=0\right\}.
    \]
    Each \(E_n\) has Lebesgue measure zero, and therefore
    \[
        E=\bigcup_{n=1}^{+\infty}E_n
    \]
    is negligible.

    On \(A_n\), we have
    \[
        x^3-n|y|<0,
    \]
    hence
    \[
        0<
        \operatorname{exp}\left(\frac{1}{x^3-n|y|}\right)
        \leq 1.
    \]
    Therefore, for \((x,y)\notin E\),
    \[
        \left|
            \frac{1}{n^2}f_n(x,y)\chi_{A_n}(x,y)
        \right|
        \leq
        \frac1{n^2}\frac{|y|}{1+x^4+y^4}.
    \]
    Since
    \[
        \sum_{n=1}^{+\infty}\frac1{n^2}<+\infty,
    \]
    the series defining \(f\) is absolutely convergent for every
    \((x,y)\notin E\). Hence \(f\) is defined almost everywhere.

    We now show that \(f\) is non-null almost everywhere. If \(y>0\), then every
    non-zero term in the series defining \(f(x,y)\) is positive. Moreover, there
    are infinitely many \(n\) such that
    \[
        x^3-n|y|<0,
    \]
    because this inequality is certainly true for all sufficiently large \(n\).
    Hence
    \[
        f(x,y)>0
        \qquad\text{for }y>0.
    \]
    Similarly, if \(y<0\), every non-zero term is negative, and therefore
    \[
        f(x,y)<0
        \qquad\text{for }y<0.
    \]
    Thus \(f(x,y)\neq0\) for every \(y\neq0\), outside the negligible set \(E\).
    Since the line \(\{y=0\}\) is negligible, \(f\) is non-null almost
    everywhere.

    Next, we prove that \(f\in\mathcal L^1(\realnumbers^2)\). From the previous
    estimate,
    \[
        |f(x,y)|
        \leq
        \left(\sum_{n=1}^{+\infty}\frac1{n^2}\right)
        \frac{|y|}{1+x^4+y^4}
    \]
    almost everywhere. It remains to observe that
    \[
        \frac{|y|}{1+x^4+y^4}\in\mathcal L^1(\realnumbers^2).
    \]
    Indeed, near the origin it is bounded, while for large \((x,y)\), writing
    \(r=\sqrt{x^2+y^2}\), one has
    \[
        1+x^4+y^4\geq C r^4,
        \qquad
        |y|\leq r,
    \]
    and therefore
    \[
        \frac{|y|}{1+x^4+y^4}\leq \frac{C}{r^3},
    \]
    which is integrable at infinity in dimension \(2\). Hence
    \[
        f\in\mathcal L^1(\realnumbers^2).
    \]

    Since the series is absolutely dominated by an integrable function, we can
    exchange sum and integral. For each \(n\), the set \(A_n\) is symmetric with
    respect to \(y\mapsto -y\). Moreover,
    \[
        \frac{y}{1+x^4+y^4}
        \operatorname{exp}\left(
            \frac{1}{x^3-n|y|}
        \right)
    \]
    is odd with respect to \(y\). Therefore
    \[
        \int_{\realnumbers^2} f_n\chi_{A_n}\,d\mu=0
    \]
    for every \(n\), and consequently
    \[
        \int_{\realnumbers^2}f\,d\mu
        =
        \sum_{n=1}^{+\infty}
        \frac1{n^2}
        \int_{\realnumbers^2} f_n\chi_{A_n}\,d\mu
        =
        0.
    \]

    We now show that
    \[
        f_n\notin\mathcal L^1(\realnumbers^2\difference A_n).
    \]
    On \(\realnumbers^2\difference A_n\), the denominator
    \(x^3-n|y|\) is non-negative. Consider the subset where \(x\in[1,2]\) and
    \(y>0\). For \(0<x^3-ny<\varepsilon\), we have
    \[
        y\sim \frac{x^3}{n},
    \]
    and the factor
    \[
        \frac{y}{1+x^4+y^4}
    \]
    is bounded from below by a positive constant depending on \(n\). Hence, on
    this region,
    \[
        |f_n(x,y)|
        \geq
        c_n
        \operatorname{exp}\left(\frac{1}{x^3-ny}\right).
    \]
    Therefore
    \[
        \int_{\realnumbers^2\difference A_n}|f_n|\,d\mu
        \geq
        c_n
        \int_1^2
        \int_{\frac{x^3-\varepsilon}{n}}^{\frac{x^3}{n}}
        \operatorname{exp}\left(\frac{1}{x^3-ny}\right)
        \,dy\,dx.
    \]
    With the change of variables
    \[
        s=x^3-ny,
        \qquad
        dy=-\frac1n\,ds,
    \]
    the inner integral becomes
    \[
        \frac1n\int_0^\varepsilon \operatorname{exp}\left(\frac1s\right)\,ds.
    \]
    This integral is infinite, because
    \[
        \operatorname{exp}\left(\frac1s\right)\geq \frac1{s^2}
    \]
    for \(s>0\) sufficiently small. Hence
    \[
        \int_{\realnumbers^2\difference A_n}|f_n|\,d\mu=+\infty,
    \]
    and so
    \[
        f_n\notin\mathcal L^1(\realnumbers^2\difference A_n).
    \]

    Finally, let \(h\colon\realnumbers\fromto\realnumbers\) be measurable. Since
    \(f\in\mathcal L^1(\realnumbers^2)\) and \(f\) is non-null almost
    everywhere, we have
    \[
        \int_{\realnumbers^3}|h(z)f(x,y)|\,d\mu(x,y,z)
        =
        \left(
            \int_{\realnumbers}|h(z)|\,dz
        \right)
        \left(
            \int_{\realnumbers^2}|f(x,y)|\,dx\,dy
        \right)
    \]
    by Tonelli's theorem. Moreover,
    \[
        0<
        \int_{\realnumbers^2}|f(x,y)|\,dx\,dy
        <+\infty.
    \]
    Therefore
    \[
        h(z)f(x,y)\in\mathcal L^1(\realnumbers^3)
        \quad\Longleftrightarrow\quad
        h\in\mathcal L^1(\realnumbers).
    \]
\end{snippetsolution}

\begin{snippetexercise}{lebesgue-measure-ex-20}{}
    For every positive \(p \in \realnumbers\), let
    \[
        E_p = \{
            (x,y,z) \in \realnumbers^3
            \suchthat x \geq 0 \land y \geq 0 \land z \geq 0 \land (xy)^p + (yz)^p + (xz)^p \leq 1
        \}
    \]
    and also let
    \[
        E_\infty = \left\{
            (x,y,z) \in \realnumbers^3 \suchthat
            x \geq 0 \land y \geq 0 \land z \geq 0 \land
            xy \leq 1 \land yz \leq 1 \land xz \leq 1
        \right\}
    \]
    \begin{enumerate}
        \item Determine whether the \set[sets] \(E_p\) for \(0 < p \leq \infty\) are bounded;
        \item Compute the Lebesgue measure of \(E_1, E_\infty\);
        \item Show that \(E_p\) has finite measure for every \(p > 0\).
    \end{enumerate}
\end{snippetexercise}

\begin{snippetsolution}{lebesgue-measure-ex-20-sol}{}
    First of all, none of the sets \(E_p\), with \(0<p\leq\infty\), is bounded.
    Indeed, for every \(x\geq0\),
    \[
        (x,0,0)\in E_p
    \]
    for every \(0<p\leq\infty\). Hence each \(E_p\) contains an unbounded
    half-line.

    We now compute the measure of \(E_1\) and \(E_\infty\). Since the coordinate
    planes have measure zero, we can restrict ourselves to \(x,y,z>0\). Define
    \[
        u=xy,\qquad v=yz,\qquad w=xz.
    \]
    Then
    \[
        x=\sqrt{\frac{uw}{v}},
        \qquad
        y=\sqrt{\frac{uv}{w}},
        \qquad
        z=\sqrt{\frac{vw}{u}}.
    \]
    This is a \(C^1\)-diffeomorphism from \((0,+\infty)^3\) onto itself. Its
    Jacobian is
    \[
        dx\,dy\,dz
        =
        \frac{1}{2\sqrt{uvw}}\,du\,dv\,dw.
    \]
    Indeed, using logarithmic coordinates,
    \[
        \log u=\log x+\log y,\qquad
        \log v=\log y+\log z,\qquad
        \log w=\log x+\log z,
    \]
    and the determinant of this linear transformation is \(2\).

    For \(E_1\), the condition becomes
    \[
        u+v+w\leq1.
    \]
    Therefore
    \[
        \mu(E_1)
        =
        \frac12
        \integral[u+v+w\leq1]{}{
            u^{-\sfrac12}v^{-\sfrac12}w^{-\sfrac12}
        }[(u,v,w)].
    \]
    By the Dirichlet integral formula,
    \[
        \integral[u+v+w\leq1]{}{
            u^{-\sfrac12}v^{-\sfrac12}w^{-\sfrac12}
        }[(u,v,w)]
        =
        \frac{\Gamma(\sfrac12)^3\Gamma(1)}{\Gamma(\sfrac52)}
        =
        \frac{\pi^{\sfrac32}}{\frac34\sqrt\pi}
        =
        \frac{4\pi}{3}.
    \]
    Hence
    \[
        \mu(E_1)=\frac{2\pi}{3}.
    \]

    For \(E_\infty\), the conditions become
    \[
        0\leq u\leq1,
        \qquad
        0\leq v\leq1,
        \qquad
        0\leq w\leq1.
    \]
    Hence
    \[
        \mu(E_\infty)
        =
        \frac12
        \integral[0][1]{u^{-\sfrac12}}[u]
        \integral[0][1]{v^{-\sfrac12}}[v]
        \integral[0][1]{w^{-\sfrac12}}[w].
    \]
    Since
    \[
        \integral[0][1]{s^{-\sfrac12}}[s]=2,
    \]
    we get
    \[
        \mu(E_\infty)
        =
        \frac12\cdot 2^3
        =
        4.
    \]

    Finally, let \(p>0\). In the variables \(u,v,w\), the set \(E_p\) is
    described by
    \[
        u^p+v^p+w^p\leq1.
    \]
    In particular,
    \[
        0\leq u\leq1,\qquad 0\leq v\leq1,\qquad 0\leq w\leq1.
    \]
    Thus
    \[
        E_p\subseteq E_\infty.
    \]
    Since \(\mu(E_\infty)=4<+\infty\), it follows that
    \[
        \mu(E_p)<+\infty
    \]
    for every \(p>0\).

    Therefore the sets \(E_p\) are all unbounded, but they all have finite
    Lebesgue measure.
\end{snippetsolution}

\end{document}